% TEMPLATE FILE — replaced on update. Do not edit.
%
% The one rules chapter of the lab book. Everything else in this document is
% dated work, in order.

\section{How this lab book works}
\label{sec:rules}

\subsection*{What it is}

The technical record of this project: calculations as they were done, the code that came out of
them, and the figures that resulted. It is written by the agent, from the researcher's
handwritten calculations and from the code they are working on.

It is \emph{not} a paper, and it is not reorganised to read like one. Later understanding does not
get retrofitted into earlier entries. If something turns out to be wrong, a later entry says so.

The narrative record --- what was worked on in a given week, what is still open --- is a separate
document, \texttt{worklog.pdf}. Nothing in this book depends on it.

\subsection*{Order}

Strictly chronological. One file per piece of work,
\texttt{labbook/entries/YYYY-MM-DD-slug.tex}, and the include list in
\texttt{labbook/entries.tex} is generated from the directory in date order by
\texttt{build\_entries.py}. There are no topical chapters, and nothing is ever moved.

Two entries on the same day are distinguished by their slug and ordered by it.

\subsection*{What earns an entry}

\begin{itemize}
  \item a handwritten calculation was pasted in;
  \item code was written or substantially changed;
  \item a figure was produced;
  \item a result was reached, or a line of work was abandoned.
\end{itemize}

Not: a typo, a rename, a formatting pass.

An entry is \emph{complete} if the work it describes could be redone from it alone --- the
equations it rests on, the code that ran, the parameters, and what came out.

\subsection*{The shape of an entry}

Every entry begins with a header comment carrying its metadata, which
\texttt{check.py} verifies, and then \verb|\entryhead|:

\begin{verbatim}
% ENTRY-DATE: 2026-08-28
% ENTRY-TITLE: Gap scaling near the band edge
% ENTRY-SUMMARY: Transcribes the handwritten gap expansion and the code it produced.

\entryhead{2026-08-28}{Gap scaling near the band edge}
          {Where the gap expansion came from and what the code does with it.}
\end{verbatim}

Then, as applicable and in this order:

\begin{enumerate}
  \item \textbf{What was asked.} One or two sentences, in the researcher's own terms.
  \item \textbf{The calculation}, inside a \texttt{transcription} environment, reproduced as
        written --- their notation, their ordering, their conventions. Nothing tidied, nothing
        re-derived. Anything illegible is marked \verb|\unreadable{...}| and asked about.
  \item \textbf{The code}, naming the file and what it computes, and how it maps onto the
        equations above.
  \item \textbf{Figures}, included from \texttt{figures/} with \verb|\includegraphics|, each
        pointing at its caption file.
  \item \textbf{What came out}, and whether it has been checked.
  \item \textbf{Flags}, if the check found something --- see below.
\end{enumerate}

\subsection*{Conventions}

\begin{itemize}
  \item The researcher's notation wins. Where a source uses a different convention, the entry states the
        translation rather than silently converting.
  \item Every number that came from a run names the script and the commit it came from, with
        \verb|\fromcode| and \verb|\atcommit|.
  \item A statement that has not been checked is marked \verb|\unchecked|. This is not an
        embarrassment; an unmarked wrong number is.
\end{itemize}

\subsection*{Flagged errors}

If the agent believes it has found a mistake in the researcher's mathematics, it says so \emph{after}
transcribing faithfully, never by editing. The flag, its evidence, the proposed fix and
\textbf{the researcher's decision} are recorded in the entry with \verb|\flag|:

\begin{verbatim}
\flag{2026-08-28}
     {The factor of $2$ in the second line of the expansion.}
     {Dimensional analysis: the left side is an energy, the right side is not.}
     {Drop the factor; the gap scaling below changes by $\sqrt{2}$.}
     {Fixed, 2026-08-29.}
\end{verbatim}

A decision to keep the original is recorded the same way. It is part of the record, not an error
to be hidden, and a later entry may revisit it.

\subsection*{Figures}

Figures come from the researcher's own plotting code, in their own style. Every figure file has a sibling
\texttt{figures/<name>.md} explaining in plain language what it shows, what the axes are, and what
produced it. \texttt{check.py} fails on a figure without one.
